
\documentclass{article}
\usepackage{colortbl}
\usepackage{makecell}
\usepackage{multirow}
\usepackage{supertabular}

\begin{document}

\newcounter{utterance}

\centering \large Interaction Transcript for game `referencegame', experiment `letter\_grids', episode 7 with reflexion{-}Qwen3.5{-}9B.
\vspace{24pt}

{ \footnotesize  \setcounter{utterance}{1}
\setlength{\tabcolsep}{0pt}
\begin{supertabular}{c@{$\;$}|p{.15\linewidth}@{}p{.15\linewidth}p{.15\linewidth}p{.15\linewidth}p{.15\linewidth}p{.15\linewidth}}
    \# & $\;$A & \multicolumn{4}{c}{Game Master} & $\;\:$B\\
    \hline

    \theutterance \stepcounter{utterance}  
    & & \multicolumn{4}{p{0.6\linewidth}}{
        \cellcolor[rgb]{0.9,0.9,0.9}{
            \makecell[{{p{\linewidth}}}]{
                \texttt{\tiny{[P1$\langle$GM]}}
                \texttt{You are given three grids, where each of them is 5 by 5 in size.} \\
\texttt{Grids have empty cells marked with "▢" and filled cells marked with "X".} \\
\texttt{Your task is to generate a referring expression that best describes the target grid while distinguishing it from the two other distractor grids.} \\
\texttt{The first grid is the target grid, and the following two grids are the distractors.} \\
\\ 
\texttt{Target grid:} \\
\\ 
\texttt{X X X X X} \\
\texttt{X □ □ □ □} \\
\texttt{X X X X X} \\
\texttt{X □ □ □ □} \\
\texttt{X X X X X} \\
\\ 
\texttt{Distractor grid 1:} \\
\\ 
\texttt{X X X X X} \\
\texttt{X □ X □ X} \\
\texttt{X X X X X} \\
\texttt{X □ X □ X} \\
\texttt{X X X X X} \\
\\ 
\texttt{Distractor grid 2:} \\
\\ 
\texttt{X X X X X} \\
\texttt{X □ □ □ □} \\
\texttt{X □ □ □ □} \\
\texttt{X □ □ □ □} \\
\texttt{X X X X X} \\
\\ 
\texttt{Instruction: Describe the target grid.} \\
\texttt{Generate the referring expression starting with the tag "Expression: " for the given target grid. Omit any other text.} \\
            }
        }
    }
    & \\ \\

    \theutterance \stepcounter{utterance}  
    & \multicolumn{4}{p{0.6\linewidth}}{
        \cellcolor[rgb]{0.9,0.9,0.9}{
            \makecell[{{p{\linewidth}}}]{
                \texttt{\tiny{[P1$\rangle$GM]}}
                \texttt{Expression: The grid with a solid border of X's and two horizontal rows of empty cells in the middle, specifically the second and fourth rows, while the first, third, and fifth rows are completely filled with X's.} \\
            }
        }
    }
    & & \\ \\

    \theutterance \stepcounter{utterance}  
    & & & \multicolumn{2}{p{0.3\linewidth}}{
        \cellcolor[rgb]{0.9,0.9,0.9}{
            \makecell[{{p{\linewidth}}}]{
                \texttt{\tiny{[GM$|$GM]}}
                \texttt{Expression: The grid with a solid border of X's and two horizontal rows of empty cells in the middle, specifically the second and fourth rows, while the first, third, and fifth rows are completely filled with X's.} \\
            }
        }
    }
    & & \\ \\

    \theutterance \stepcounter{utterance}  
    & & \multicolumn{4}{p{0.6\linewidth}}{
        \cellcolor[rgb]{0.9,0.9,0.9}{
            \makecell[{{p{\linewidth}}}]{
                \texttt{\tiny{[GM$\rangle$P2]}}
                \texttt{You are given three grids, where each of them is 5 by 5 in size.} \\
\texttt{Grids have empty cells marked with "▢" and filled cells marked with "X".} \\
\texttt{You are also given a referring expression that describes one of the given grids.} \\
\texttt{Your task is to select the grid that matches the given referring expression.} \\
\texttt{Generate only the number (in text) of the grid that matches the given expression by selecting first, second, or third.} \\
\\ 
\texttt{First:} \\
\\ 
\texttt{X X X X X} \\
\texttt{X □ X □ X} \\
\texttt{X X X X X} \\
\texttt{X □ X □ X} \\
\texttt{X X X X X} \\
\\ 
\texttt{Second:} \\
\\ 
\texttt{X X X X X} \\
\texttt{X □ □ □ □} \\
\texttt{X X X X X} \\
\texttt{X □ □ □ □} \\
\texttt{X X X X X} \\
\\ 
\texttt{Third:} \\
\\ 
\texttt{X X X X X} \\
\texttt{X □ □ □ □} \\
\texttt{X □ □ □ □} \\
\texttt{X □ □ □ □} \\
\texttt{X X X X X} \\
\\ 
\texttt{Expression: The grid with a solid border of X's and two horizontal rows of empty cells in the middle, specifically the second and fourth rows, while the first, third, and fifth rows are completely filled with X's.} \\
\texttt{Question: Which grid does the expression refer to?} \\
\texttt{Start with the tag "Answer: ", followed by your selection. Omit any other text.} \\
            }
        }
    }
    & \\ \\

    \theutterance \stepcounter{utterance}  
    & & & \multicolumn{4}{p{0.6\linewidth}}{
        \cellcolor[rgb]{0.9,0.9,0.9}{
            \makecell[{{p{\linewidth}}}]{
                \texttt{\tiny{[GM$\langle$P2]}}
                \texttt{Answer: First} \\
            }
        }
    }
     \\ \\

    \theutterance \stepcounter{utterance}  
    & & & \multicolumn{2}{p{0.3\linewidth}}{
        \cellcolor[rgb]{0.9,0.9,0.9}{
            \makecell[{{p{\linewidth}}}]{
                \texttt{\tiny{[GM$|$GM]}}
                \texttt{first} \\
            }
        }
    }
    & & \\ \\

    \theutterance \stepcounter{utterance}  
    & \multicolumn{4}{p{0.6\linewidth}}{
        \cellcolor[rgb]{0.95,0.95,0.95}{
            \makecell[{{p{\linewidth}}}]{
                \texttt{\tiny{[P2$\rightarrow$P1]}}
                \texttt{Post{-}episode Reflexion (not sent to GM):} \\
\texttt{New plan: My previous description was too vague by using relative terms like "middle" and "specifically the second and fourth rows" without explicitly stating the row indices, which likely caused ambiguity with the distractors. Next time, I will construct the expression by explicitly defining the grid dimensions (5x5), specifying the exact row numbers (2 and 4) that contain the empty cells, and contrasting them directly with the fully filled rows (1, 3, and 5) to ensure a unique and unambiguous distinction from the distractor grids.} \\
            }
        }
    }
    & & \\ \\

    \theutterance \stepcounter{utterance}  
    & & & \multicolumn{4}{p{0.6\linewidth}}{
        \cellcolor[rgb]{0.95,0.95,0.95}{
            \makecell[{{p{\linewidth}}}]{
                \texttt{\tiny{[P1$\rightarrow$P2]}}
                \texttt{Post{-}episode Reflexion (not sent to GM):} \\
\texttt{New plan: I incorrectly identified the "First" grid because I focused on the alternating pattern of empty cells in rows 2 and 4, but I failed to verify the specific row indices mentioned in the expression. The expression explicitly requires empty cells in the **second and fourth** rows, whereas the First grid has empty cells in the second and fourth rows but the description of "two horizontal rows... in the middle" combined with the visual check reveals the First grid actually has empty cells in rows 2 and 4? Wait, let me re{-}examine the First grid: Row 2 is `X □ X □ X` (empty at cols 2,4) and Row 4 is `X □ X □ X` (empty at cols 2,4). The Second grid has Row 2 as `X □ □ □ □` (empty at cols 2,3,4,5) and Row 4 as `X □ □ □ □`. The expression says "two horizontal rows of empty cells", implying the cells within those rows are empty, or perhaps it means the rows themselves are empty except for the border? Re{-}reading: "two horizontal rows of empty cells in the middle". The First grid has rows with *some* empty cells, not rows *of* empty cells. The Second grid has rows 2 and 4 where the inner cells are empty, forming a block. Actually, looking closer at the First grid: Row 2 is `X □ X □ X`. This is not a row of empty cells. The Second grid Row 2 is `X □ □ □ □`. This is a row with mostly empty cells. The Third grid has three such rows. The expression says "two horizontal rows". The First grid has two rows with empty cells (2 and 4), but they are not "rows of empty cells" in the sense of being filled with emptiness, they are alternating. However, the most likely error is misinterpreting "two horizontal rows of empty cells" as "rows containing empty cells" vs "rows that are empty". Let's look at the First grid again. Row 2: X, empty, X, empty, X. Row 4: X, empty, X, empty, X. This fits "two rows... specifically the second and fourth". But wait, the Second grid Row 2 is X, empty, empty, empty, empty. That is a row of empty cells (mostly). The First grid is alternating. The expression says "two horizontal rows of empty cells". This phrasing usually implies the cells in those rows are empty. The First grid does not have rows of empty cells; it has rows with alternating X and empty. The Second grid has rows 2 and 4 that are `X □ □ □ □`, which are rows of empty cells (with a border). The First grid is `X □ X □ X`. The key distinction is} \\
            }
        }
    }
     \\ \\

\end{supertabular}
}

\end{document}
